%=====================================================================
\chapter{Resources}\label{app:resources}
%=====================================================================
\normalfigures

\section{Textbooks}

\rowcolors{2}{white}{lightbg}
\begin{longtable}{@{}L{6.6cm}L{8.4cm}@{}}
\toprule
\rowcolor{primary!15}
\textbf{Book} & \textbf{Use it for} \\
\midrule
\endhead
Tom M.\ Mitchell, \emph{Machine Learning}, McGraw-Hill, 1997 (HEC set text) &
Chapters 1, 4, 9, 10, 11, 13 of this handout follow its treatment; its EnjoySport and
PlayTennis tables are the release and sprint data of Chapters 4 and 9, renamed \\
Kevin P.\ Murphy, \emph{Machine Learning: A Probabilistic Perspective}, MIT Press, 2012
(HEC set text) & The probabilistic view: Bayesian learning, mixtures and EM, HMMs \\
Kevin P.\ Murphy, \emph{Probabilistic Machine Learning: An Introduction}, MIT Press,
2022 & The updated successor to the set text; freely available from the author \\
Gareth James et al., \emph{An Introduction to Statistical Learning with Applications in
Python}, Springer, 2023 & The gentlest rigorous treatment of regression, classification,
resampling and trees; free online \\
Aur\'elien G\'eron, \emph{Hands-On Machine Learning with Scikit-Learn, Keras and
TensorFlow}, 3rd ed., O'Reilly, 2022 & Practical scikit-learn; the labs' style \\
Richard S.\ Sutton and Andrew G.\ Barto, \emph{Reinforcement Learning: An Introduction},
2nd ed., MIT Press, 2018 & Chapters 19--20; bandits, Monte Carlo, TD, Q-learning; free
online \\
Christopher M.\ Bishop, \emph{Pattern Recognition and Machine Learning}, Springer, 2006
& A deeper mathematical treatment of mixtures, kernels and neural networks \\
\bottomrule
\end{longtable}

\section{Online Courses}

HEC 2025's AI specialization recommends online courses as companions to its core
courses, and for \emph{Machine Learning} names the Machine Learning Specialization
(DeepLearning.AI and Stanford, on Coursera) and Columbia University's Machine Learning
course on edX. Both align well with Parts II and III of this handout. Stanford's CS229
lecture notes are a free, more mathematical supplement.

\section{Datasets}

\rowcolors{2}{white}{lightbg}
\begin{longtable}{@{}L{4.6cm}L{10.4cm}@{}}
\toprule
\rowcolor{primary!15}
\textbf{Source} & \textbf{Notes} \\
\midrule
\endhead
Synthetic project data & Every lab builds its own tasks, projects, tickets or code
changes with NumPy or scikit-learn's \texttt{make\_*} generators, so it runs offline
and its answers are known \\
Open-source issue trackers & The Apache Software Foundation's public Jira and GitHub
issues: real tickets with text, history, status changes and resolution times \\
PROMISE repository & Public software-engineering datasets, among them effort-estimation
and defect-prediction data \\
UCI, OpenML, Kaggle & General collections; search them for software-project,
effort and defect data, and read other people's notebooks critically \\
Your employer's tracker & Often the most meaningful data; requires written permission
and the removal of personal identifiers (\chref{ethics}) \\
\bottomrule
\end{longtable}

\section{Tools}

Anaconda (Python distribution), Jupyter notebooks for exploration, a plain editor or
VS Code for scripts, Git for version control, and the scikit-learn user guide --- the
single most useful reference for the practical side of this course.

\section{The Summer 2026 Handout}

The department's \emph{Machine Learning --- Core Concepts} handout (Summer 2026) is a
24-page overview written for an eight-week summer session. It covers the supervised,
unsupervised and neural-network material of Parts II--IV in outline, but not concept
learning, version spaces, self-organizing maps, semi-supervised EM, hidden Markov models
or Markov decision processes, which the HEC outline requires. Use it as a revision
summary alongside this volume, not in place of it.
